% Options for packages loaded elsewhere
% Options for packages loaded elsewhere
\PassOptionsToPackage{unicode}{hyperref}
\PassOptionsToPackage{hyphens}{url}
\PassOptionsToPackage{dvipsnames,svgnames,x11names}{xcolor}
%
\documentclass[
  letterpaper,
  DIV=11,
  numbers=noendperiod]{scrartcl}
\usepackage{xcolor}
\usepackage{amsmath,amssymb}
\setcounter{secnumdepth}{-\maxdimen} % remove section numbering
\usepackage{iftex}
\ifPDFTeX
  \usepackage[T1]{fontenc}
  \usepackage[utf8]{inputenc}
  \usepackage{textcomp} % provide euro and other symbols
\else % if luatex or xetex
  \usepackage{unicode-math} % this also loads fontspec
  \defaultfontfeatures{Scale=MatchLowercase}
  \defaultfontfeatures[\rmfamily]{Ligatures=TeX,Scale=1}
\fi
\usepackage{lmodern}
\ifPDFTeX\else
  % xetex/luatex font selection
\fi
% Use upquote if available, for straight quotes in verbatim environments
\IfFileExists{upquote.sty}{\usepackage{upquote}}{}
\IfFileExists{microtype.sty}{% use microtype if available
  \usepackage[]{microtype}
  \UseMicrotypeSet[protrusion]{basicmath} % disable protrusion for tt fonts
}{}
\makeatletter
\@ifundefined{KOMAClassName}{% if non-KOMA class
  \IfFileExists{parskip.sty}{%
    \usepackage{parskip}
  }{% else
    \setlength{\parindent}{0pt}
    \setlength{\parskip}{6pt plus 2pt minus 1pt}}
}{% if KOMA class
  \KOMAoptions{parskip=half}}
\makeatother
% Make \paragraph and \subparagraph free-standing
\makeatletter
\ifx\paragraph\undefined\else
  \let\oldparagraph\paragraph
  \renewcommand{\paragraph}{
    \@ifstar
      \xxxParagraphStar
      \xxxParagraphNoStar
  }
  \newcommand{\xxxParagraphStar}[1]{\oldparagraph*{#1}\mbox{}}
  \newcommand{\xxxParagraphNoStar}[1]{\oldparagraph{#1}\mbox{}}
\fi
\ifx\subparagraph\undefined\else
  \let\oldsubparagraph\subparagraph
  \renewcommand{\subparagraph}{
    \@ifstar
      \xxxSubParagraphStar
      \xxxSubParagraphNoStar
  }
  \newcommand{\xxxSubParagraphStar}[1]{\oldsubparagraph*{#1}\mbox{}}
  \newcommand{\xxxSubParagraphNoStar}[1]{\oldsubparagraph{#1}\mbox{}}
\fi
\makeatother


\usepackage{longtable,booktabs,array}
\newcounter{none} % for unnumbered tables
\usepackage{calc} % for calculating minipage widths
% Correct order of tables after \paragraph or \subparagraph
\usepackage{etoolbox}
\makeatletter
\patchcmd\longtable{\par}{\if@noskipsec\mbox{}\fi\par}{}{}
\makeatother
% Allow footnotes in longtable head/foot
\IfFileExists{footnotehyper.sty}{\usepackage{footnotehyper}}{\usepackage{footnote}}
\makesavenoteenv{longtable}
\usepackage{graphicx}
\makeatletter
\newsavebox\pandoc@box
\newcommand*\pandocbounded[1]{% scales image to fit in text height/width
  \sbox\pandoc@box{#1}%
  \Gscale@div\@tempa{\textheight}{\dimexpr\ht\pandoc@box+\dp\pandoc@box\relax}%
  \Gscale@div\@tempb{\linewidth}{\wd\pandoc@box}%
  \ifdim\@tempb\p@<\@tempa\p@\let\@tempa\@tempb\fi% select the smaller of both
  \ifdim\@tempa\p@<\p@\scalebox{\@tempa}{\usebox\pandoc@box}%
  \else\usebox{\pandoc@box}%
  \fi%
}
% Set default figure placement to htbp
\def\fps@figure{htbp}
\makeatother





\setlength{\emergencystretch}{3em} % prevent overfull lines

\providecommand{\tightlist}{%
  \setlength{\itemsep}{0pt}\setlength{\parskip}{0pt}}



 


\KOMAoption{captions}{tableheading}
\makeatletter
\@ifpackageloaded{caption}{}{\usepackage{caption}}
\AtBeginDocument{%
\ifdefined\contentsname
  \renewcommand*\contentsname{Table of contents}
\else
  \newcommand\contentsname{Table of contents}
\fi
\ifdefined\listfigurename
  \renewcommand*\listfigurename{List of Figures}
\else
  \newcommand\listfigurename{List of Figures}
\fi
\ifdefined\listtablename
  \renewcommand*\listtablename{List of Tables}
\else
  \newcommand\listtablename{List of Tables}
\fi
\ifdefined\figurename
  \renewcommand*\figurename{Figure}
\else
  \newcommand\figurename{Figure}
\fi
\ifdefined\tablename
  \renewcommand*\tablename{Table}
\else
  \newcommand\tablename{Table}
\fi
}
\@ifpackageloaded{float}{}{\usepackage{float}}
\floatstyle{ruled}
\@ifundefined{c@chapter}{\newfloat{codelisting}{h}{lop}}{\newfloat{codelisting}{h}{lop}[chapter]}
\floatname{codelisting}{Listing}
\newcommand*\listoflistings{\listof{codelisting}{List of Listings}}
\makeatother
\makeatletter
\makeatother
\makeatletter
\@ifpackageloaded{caption}{}{\usepackage{caption}}
\@ifpackageloaded{subcaption}{}{\usepackage{subcaption}}
\makeatother
\usepackage{bookmark}
\IfFileExists{xurl.sty}{\usepackage{xurl}}{} % add URL line breaks if available
\urlstyle{same}
\hypersetup{
  pdftitle={Guide to the Interactive HTML Outputs},
  colorlinks=true,
  linkcolor={blue},
  filecolor={Maroon},
  citecolor={Blue},
  urlcolor={Blue},
  pdfcreator={LaTeX via pandoc}}


\title{Guide to the Interactive HTML Outputs}
\author{}
\date{June 9, 2026}
\begin{document}
\maketitle


\subsection{Overview}\label{overview}

This folder contains a set of interactive HTML files, all generated by
the analysis scripts included in the online repository, designed to
allow readers to visually inspect the primary analytical pipeline of
this study. Each file is a standalone web page: open it in any modern
browser by double-clicking, no server, internet connection, or
additional software is required.

All interactive files live under \texttt{analysis/output/html/}:

\begin{verbatim}
analysis/output/html/
├── scar_alignment_svd.html              # technological alignment diagnostic
├── scar_alignment_lin2024.html          # morphological alignment diagnostic
├── kde_sphere_interactive/              # scar vector density on a sphere
│   ├── kde_sphere_EXP.html
│   ├── kde_sphere_IM.html
│   └── kde_sphere_SDG.html
└── reconstruction_interactive/          # spherical harmonic reconstructions
    ├── interactive_EXP.html
    ├── interactive_IM.html
    ├── interactive_SDG.html
    ├── interactive_EXP_CoIA_axes.html
    └── interactive_SDG_CoIA_axes.html
\end{verbatim}

Three dataset codes appear throughout the filenames:

\begin{itemize}
\tightlist
\item
  \textbf{IM:} the idealized digital core models.
\item
  \textbf{EXP:} the experimental knapped cores.
\item
  \textbf{SDG:} the archaeological Sandinggai cores.
\end{itemize}

The sections below describe what each file shows and how to read it,
together with the script that generates it.

\subsection{Alignment diagnostics}\label{alignment-diagnostics}

These two pages visualise, specimen by specimen, the orientation
standardisation that is applied to the raw flaking scar vectors before
any SPHARM analysis. In both pages a dropdown at the top selects the
specimen, every panel is a fully rotatable/zoomable 3D plot, and a short
caption beneath each panel states what that step does.

\subsubsection{\texorpdfstring{\texttt{scar\_alignment\_svd.html}}{scar\_alignment\_svd.html}}\label{scar_alignment_svd.html}

Generated by \texttt{analysis/scripts/r\_alignment/align\_svd.R}. This
is the alignment used for the main analysis. It shows the four step SVD
pipeline as four side-by-side panels:

\begin{itemize}
\tightlist
\item
  \textbf{Step 0: Raw data.} Scar start/end points in their original
  coordinates, with the SVD-fitted best-fit plane and its normal.
\item
  \textbf{Step 1: Rotate.} The best-fit plane normal (the third right
  singular vector of the unit scar direction vectors) is rotated onto
  the Z-axis, so the plane lies flat in XY.
\item
  \textbf{Step 2: Translate.} The centroid of all scar endpoints is
  moved to the origin, removing positional bias between specimens.
\item
  \textbf{Step 3: Rotate XY.} The dominant in-plane scar direction
  (first singular vector of the XY projections) is rotated onto the
  X-axis, so every specimen shares a common reference frame.
\end{itemize}

\subsubsection{\texorpdfstring{\texttt{scar\_alignment\_lin2024.html}}{scar\_alignment\_lin2024.html}}\label{scar_alignment_lin2024.html}

Generated by \texttt{analysis/scripts/r\_alignment/align\_lin2024.R}.
This page documents the alternative alignment of Lin et al., 2024, used
in the rotational-invariance validation rather than the main pipeline.
It has three panels:

\begin{itemize}
\tightlist
\item
  \textbf{Step 0: Direction vectors / reference normal.} Unit direction
  vectors are computed from scar endpoints; the morphological plane
  normal (\texttt{Norm\_X/Y/Z}, measured externally in Geomagic Wrap
  software) serves as the reference orientation.
\item
  \textbf{Step 1: Rotate.} A Rodrigues rotation aligns the measured
  morphological normal with the Z-axis, placing the best-fit plane in
  XY.
\item
  \textbf{Step 2: Translate.} The longest scar is identified and its
  rotated start-point is moved to (0, 0), anchoring the origin while
  preserving the absolute spatial relationships between scars.
\end{itemize}

\subsection{Scar vector density on the
sphere}\label{scar-vector-density-on-the-sphere}

\subsubsection{\texorpdfstring{\texttt{kde\_sphere\_interactive/kde\_sphere\_\{EXP,IM,SDG\}.html}}{kde\_sphere\_interactive/kde\_sphere\_\{EXP,IM,SDG\}.html}}\label{kde_sphere_interactivekde_sphere_expimsdg.html}

Generated by
\texttt{analysis/scripts/SPHARM\_modules/spherical\_kde\_interactive\_export.py},
one file per dataset. Each page renders the von Mises--Fisher kernel
density estimate (KDE) of a specimen's flaking-scar direction vectors on
a single unit sphere, coloured by density. This is the spherical density
that the SP-SPHARM descriptors are computed from, shown here as a
directly inspectable surface.

Only the raw direction vectors are embedded; the KDE itself is
recomputed in the browser, so the density responds live to the bandwidth
control.

Controls:

\begin{itemize}
\tightlist
\item
  \textbf{Specimen dropdown}, grouped by core type.
\item
  \textbf{Bandwidth slider (\texttt{bw}):} recomputes the KDE on the
  fly; the concentration parameter is \emph{κ = 1/bw²}.
\item
  \textbf{Type Mean:} pools the direction vectors of every specimen of
  the selected type and re-fits the KDE, giving a per-type average
  pattern.
\item
  \textbf{Scars:} toggles the scatter of individual scar vector points
  over the density surface.
\item
  \textbf{Contours:} overlays iso-density contour lines.
\item
  \textbf{2D Map:} an equirectangular density map.
\item
  \textbf{Graticule}, \textbf{view presets} (Iso / Top / Front),
  \textbf{material} selector, \textbf{Info} panel, a density
  \textbf{colour bar}, and \textbf{PNG} export.
\end{itemize}

\subsection{SPHARM reconstructions}\label{spharm-reconstructions}

\subsubsection{\texorpdfstring{\texttt{reconstruction\_interactive/interactive\_\{EXP,IM,SDG\}.html}}{reconstruction\_interactive/interactive\_\{EXP,IM,SDG\}.html}}\label{reconstruction_interactiveinteractive_expimsdg.html}

Generated by
\texttt{analysis/scripts/SPHARM\_modules/spharm\_interactive\_export.py},
one file per dataset. Each page is a dual viewport: the Morphology
reconstruction (M-SPHARM, left) and the Scar pattern reconstruction
(SP-SPHARM, right) for the same specimen, with synchronised rotation and
zoom. The surfaces are synthesised in the browser in real time from the
stored spherical-harmonic coefficients (up to degree \emph{l} = 20).

Controls:

\begin{itemize}
\tightlist
\item
  \textbf{Specimen dropdown}, grouped by core type.
\item
  \textbf{Degree slider (\texttt{l\ ≤})} with a \textbf{play button:}
  rebuilds the surface using only harmonics up to the chosen degree, and
  animates \emph{l} = 1 → 20 to show how detail accumulates as higher
  frequencies are added.
\item
  \textbf{Type Mean:} averages the coefficients across all specimens of
  the selected type and reconstructs that mean shape.
\item
  \textbf{Colormap:} shades the surface by radial deviation from its
  mean radius (blue = inward, red = outward), with a colour bar.
\item
  \textbf{Wire} (wireframe), \textbf{material} selector, \textbf{view
  presets}, \textbf{Info} panel, plus \textbf{PNG} screenshot and
  \textbf{OBJ} mesh export.
\end{itemize}

\subsubsection{\texorpdfstring{\texttt{reconstruction\_interactive/interactive\_\{EXP,SDG\}\_CoIA\_axes.html}}{reconstruction\_interactive/interactive\_\{EXP,SDG\}\_CoIA\_axes.html}}\label{reconstruction_interactiveinteractive_expsdg_coia_axes.html}

Generated by the \texttt{export\_axis\_interactive\_html()} routine in
\texttt{analysis/scripts/SPHARM\_modules/spharm\_reconstruction.py}
(reusing the viewer above). These pages exist only for the two datasets
that have a co-inertia analysis (CoIA): \textbf{EXP} and \textbf{SDG}.

Instead of individual specimens, the dropdown steps through synthetic
average shapes sampled along the CoIA axes that summarise the
covariation between the morphology block and the scar pattern block. For
each of \emph{Morph Axis 1/2} and \emph{Scar Axis 1/2}, the axis is
sampled at a sequence of positions (across a robust percentile range);
at each position a Gaussian-kernel-weighted mean of the coefficients is
reconstructed. The Info panel reports the axis, the axis coordinate of
the current step, the kernel bandwidth, and the effective sample size
behind each averaged shape.




\end{document}
